% exercises/solutions-01.tex — Full solutions for Chapter 1
% Rendered in the instructor edition only (guarded in the chapter).

% ---------- Exercise 1.1 -------------------------------------------------
\begin{solution}{ex:1-1}
Apply AM--GM (\cref{thm:amgm2}) to the pair
$\big(x, \tfrac{4}{x}\big)$, whose product is $4$:
\[
  x + \frac{4}{x} \;\ge\; 2\sqrt{x \cdot \frac{4}{x}} = 2\sqrt4 = 4 .
\]
Equality holds iff $x = \tfrac{4}{x}$, i.e.\ $x = 2$.
\answer{$x + 4/x = 4$ exactly at $x = 2$.}
\end{solution}

% ---------- Exercise 1.2 -------------------------------------------------
\begin{solution}{ex:1-2}
Expanding,
\[
  (a+b)\left(\frac1a + \frac1b\right)
  = 1 + \frac{a}{b} + \frac{b}{a} + 1
  = 2 + \frac{a}{b} + \frac{b}{a}.
\]
By AM--GM on the pair $\big(\tfrac{a}{b}, \tfrac{b}{a}\big)$, whose
product is $1$,
\[
  \frac{a}{b} + \frac{b}{a} \ge 2 ,
\]
so the original expression is at least $4$. Equality holds iff
$\tfrac{a}{b} = \tfrac{b}{a}$, i.e.\ $a = b$.
\end{solution}

% ---------- Exercise 1.3 -------------------------------------------------
\begin{solution}{ex:1-3}
By Cauchy--Schwarz (\cref{thm:cs}) applied to $(1,1,1)$ and
$(a,b,c)$,
\[
  (a+b+c)^2 \;\le\; 3\,(a^2+b^2+c^2).
\]
Since $a+b+c = 1$, this gives $a^2+b^2+c^2 \ge \tfrac13$. Equality
holds iff $(a,b,c)$ is proportional to $(1,1,1)$, i.e.\
$a=b=c=\tfrac13$.
\end{solution}

% ---------- Exercise 1.4 -------------------------------------------------
\begin{solution}{ex:1-4}
Multiply the target by $2$ and rearrange:
\[
  2(a^2+b^2+c^2) - 2(ab+bc+ca)
  = (a-b)^2 + (b-c)^2 + (c-a)^2 \;\ge\; 0,
\]
a sum of squares. Equality holds iff $a = b = c$. Note this is
exactly the identity used in \cref{ex:nesbitt}.
\end{solution}

% ---------- Exercise 1.5 -------------------------------------------------
\begin{solution}{ex:1-5}
By AM--GM for three variables (\cref{thm:amgm}, $n=3$),
\[
  \frac{x+y+z}{3} \;\ge\; \sqrt[3]{xyz} = \sqrt[3]{1} = 1,
\]
so $x+y+z \ge 3$. Equality holds iff $x=y=z=1$.
\end{solution}

% ---------- Exercise 1.6 -------------------------------------------------
\begin{solution}{ex:1-6}
Apply \cref{thm:amgm2} to each factor:
\[
  x+y \ge 2\sqrt{xy}, \qquad
  y+z \ge 2\sqrt{yz}, \qquad
  z+x \ge 2\sqrt{zx}.
\]
All three left sides are positive, so multiplying the inequalities is
legitimate:
\[
  (x+y)(y+z)(z+x) \;\ge\; 8\sqrt{xy \cdot yz \cdot zx}
  = 8\sqrt{x^2y^2z^2} = 8xyz .
\]
Equality holds iff $x=y=z$.
\end{solution}

% ---------- Exercise 1.7 -------------------------------------------------
\begin{solution}{ex:1-7}
By Cauchy--Schwarz applied to
$(\sqrt a, \sqrt b, \sqrt c)$ and $(1,1,1)$,
\[
  (\sqrt a + \sqrt b + \sqrt c)^2
  \;\le\; 3\,(a+b+c) = 18,
\]
so $\sqrt a + \sqrt b + \sqrt c \le 3\sqrt2$.
\answer{Maximum $= 3\sqrt2$, attained only at $(a,b,c) = (2,2,2)$.}
Equality in Cauchy--Schwarz requires
$(\sqrt a, \sqrt b, \sqrt c)$ proportional to $(1,1,1)$, i.e.\
$a=b=c$; combined with $a+b+c=6$ this forces $a=b=c=2$, which
indeed gives $3\sqrt2$. Since the domain $\{a,b,c \ge 0,\;
a+b+c=6\}$ is compact and the objective is continuous, the maximum
exists, so the bound is sharp.
\end{solution}

% ---------- Exercise 1.8 -------------------------------------------------
\begin{solution}{ex:1-8}
Apply the Engel form (\cref{cor:engel}) to the three fractions:
\[
  \frac{x^2}{y+z} + \frac{y^2}{z+x} + \frac{z^2}{x+y}
  \;\ge\;
  \frac{(x+y+z)^2}{(y+z)+(z+x)+(x+y)}
  = \frac{(x+y+z)^2}{2(x+y+z)}
  = \frac{x+y+z}{2} .
\]
Equality holds iff $\tfrac{x}{y+z} = \tfrac{y}{z+x} = \tfrac{z}{x+y}$,
which (cross-multiplying and subtracting) forces $x=y=z$.
\end{solution}

% ---------- Challenge 1 ---------------------------------------------------
\begin{solution}{ch:1-1}
Following the hint, we prove the pointwise bound
\[
  \frac{a^3}{a^2+ab+b^2} \;\ge\; \frac{2a-b}{3} .
\]
Since $a^2+ab+b^2 > 0$, clearing denominators shows the difference is
\[
  3a^3 - (2a-b)(a^2+ab+b^2)
  \;=\; 3a^3 - \big(2a^3 + a^2b + ab^2 - b^3\big)
  \;=\; (a+b)(a-b)^2 \;\ge\; 0,
\]
where the factorization is verified by expanding
$(a+b)(a-b)^2 = a^3 - a^2b - ab^2 + b^3$. Equality holds iff $a=b$.
Applying the same bound to the other two cyclic terms and summing,
\[
  \sum_{\text{cyc}} \frac{a^3}{a^2+ab+b^2}
  \;\ge\;
  \frac{(2a-b) + (2b-c) + (2c-a)}{3}
  = \frac{a+b+c}{3} .
\]
Equality requires $a=b$, $b=c$, and $c=a$ simultaneously---that is,
$a=b=c$, where both sides indeed agree (each fraction equals
$\tfrac{a}{3}$).
\end{solution}

% ---------- Challenge 2 (IMO 2001/2) --------------------------------------
\begin{solution}{ch:1-2}
Write $S = a+b+c$. We chain three estimates.

\emph{Step 1 (Engel).} Writing each term as a square over a product,
\[
  \sum_{\text{cyc}} \frac{a}{\sqrt{a^2+8bc}}
  = \sum_{\text{cyc}} \frac{a^2}{a\sqrt{a^2+8bc}}
  \;\ge\; \frac{S^2}{\displaystyle\sum_{\text{cyc}} a\sqrt{a^2+8bc}}
  \;\eqqcolon\; \frac{S^2}{T},
\]
by \cref{cor:engel} applied with numerators $a$ and denominators
$b = a\sqrt{a^2+8bc}$.

\emph{Step 2 (Cauchy--Schwarz).} Apply \cref{thm:cs} to the vectors
$\big(\sqrt a\big)$ and $\big(\sqrt{a\,(a^2+8bc)}\big)$:
\[
  T^2
  = \left(\sum_{\text{cyc}} \sqrt{a} \cdot
      \sqrt{a(a^2+8bc)}\right)^{\!2}
  \le S \cdot \left(a^3+b^3+c^3 + 24abc\right).
\]

\emph{Step 3 (expansion + AM--GM).} We claim
\[
  a^3+b^3+c^3+24abc \;\le\; S^3 .
\]
Expanding,
\[
  S^3 = a^3+b^3+c^3
  + 3\!\!\sum_{\text{sym}}\! a^2b \; + 6abc,
\]
where $\sum_{\text{sym}} a^2b = a^2b + a^2c + ab^2 + b^2c + ac^2 +
bc^2$. It remains to show $18abc \le 3\sum_{\text{sym}} a^2b$, i.e.\
$6abc \le \sum_{\text{sym}} a^2b$. This is AM--GM
(\cref{thm:amgm}) applied to the \emph{six} terms of the symmetric
sum: their product is
\[
  (a^2b)(a^2c)(ab^2)(b^2c)(ac^2)(bc^2) = a^6 b^6 c^6,
\]
so their geometric mean is $abc$ and their sum is at least $6abc$.

\emph{Conclusion.} Combining,
\[
  T \le \sqrt{S \cdot S^3} = S^2
  \qquad\Longrightarrow\qquad
  \sum_{\text{cyc}} \frac{a}{\sqrt{a^2+8bc}}
  \;\ge\; \frac{S^2}{T} \ge \frac{S^2}{S^2} = 1 .
\]
Equality in every step requires $a=b=c$: then each term equals
$\tfrac{a}{\sqrt{9a^2}} = \tfrac13$ and the sum is exactly $1$.
\end{solution}
